\section*{Bab \ref{ch:g6:decimals} --- Bilangan Desimal}

\begin{solution}{exo:g6:decimals:1}
$5.37$; \quad $0.9$; \quad $12.004$; \quad $8.05$.
\end{solution}

\begin{solution}{exo:g6:decimals:2}
Angka persepuluhan: $3$; angka perseratusan: $5$; angka $8$ membilang puluhan.
\[
86.354 = 8 \times 10 + 6 + 3 \times \frac{1}{10} + 5 \times
\frac{1}{100} + 4 \times \frac{1}{1000}.
\]
\end{solution}

\begin{solution}{exo:g6:decimals:3}
$5.3 > 5.29$ (bandingkan $5.30$ dengan $5.29$); \quad
$0.7 = 0.70$ (\omterm{def:g1:counting:zero}{nol} di ujung tidak mengubah apa pun); \quad
$2.09 < 2.9$ (persepuluhan: $0 < 9$); \quad
$14.5 > 14.49$.
\end{solution}

\begin{solution}{exo:g6:decimals:4}
Lengkapi sampai dua angka desimal: $4.20$; $4.05$; $4.51$; $4.15$; $4.50$. Urutannya:
\[
4.05 < 4.15 < 4.2 < 4.5 < 4.51 .
\]
\end{solution}

\begin{solution}{exo:g6:decimals:5}
Setiap tanda bernilai satu persepuluhan. $A$: $6.3$; \quad $B$: $6.7$; \quad
$C$: $7.2$.
\end{solution}

\begin{solution}{exo:g6:decimals:6}
$45.80 + 7.65 = 53.45$.

$23.40 - 8.72 = 14.68$ (periksa: $14.68 + 8.72 = 23.4$).

$56 \times 24 = 1344$, dan ada dua angka desimal pada faktornya:
$5.6 \times 2.4 = 13.44$.
\end{solution}

\begin{solution}{exo:g6:decimals:7}
$6.42 \times 10 = 64.2$; \quad
$0.35 \times 1000 = 350$; \quad
$78.1 \div 100 = 0.781$; \quad
$5 \div 1000 = 0.005$.
\end{solution}

\begin{solution}{exo:g6:decimals:8}
$2.4$ m $= 240$ cm; \quad $370$ g $= 0.37$ kg (bagilah dengan $1000$);
\quad $0.85$ km $= 850$ m; \quad $56$ mm $= 0.056$ m.
\end{solution}

\begin{solution}{exo:g6:decimals:9}
Ke satuan: $24$ (angka berikutnya $8$: \omterm{def:g4:numbers:round}{bulatkan} ke atas). Ke
persepuluhan: $23.9$. Ke perseratusan: $23.87$. Dan $9.97$ ke
persepuluhan: angka berikutnya $7$, jadi \omterm{def:g4:numbers:round}{bulatkan} naik $9$ persepuluhannya --- $10.0$.
\end{solution}

\begin{solution}{exo:g6:decimals:10}
Harga ketiga baguette itu: $3 \times 1.15 = 3.45$. Kembaliannya:
$5 - 3.45 = 1.55$.
\end{solution}

\begin{solution}{exo:g6:decimals:11}
Antara $7.4$ dan $7.5$: misalnya $7.45$ (atau $7.41$, $7.499$,
\dots). Antara $3.99$ dan $4$: misalnya $3.995$. Pada kedua keadaan
itu ada \emph{tak hingga banyak} bilangan semacam itu: kita selalu
dapat menambahkan tempat desimal lagi.
\end{solution}

\begin{solution}{exo:g6:decimals:12}
Terbesar: letakkan angka terbesar di depan dan titiknya selambat
mungkin: $85.2$. Terkecil: angka terkecil di depan dan titiknya
sedini mungkin: $2.58$.
\end{solution}

\begin{solution}{pb:g6:decimals:1}
\textbf{1.} $3$ yang lebih besar: setelah dilengkapi, $3.000 > 2.999$.
\omterm{ex:g1:subtraction:difference}{Selisihnya} $3.000 - 2.999 = 0.001$, yaitu satu perseribu.

\textbf{2.} Tepat di antara $7.4 = 7.40$ dan $7.5 = 7.50$ terletak
\[
7.41,\ 7.42,\ 7.43,\ 7.44,\ 7.45,\ 7.46,\ 7.47,\ 7.48,\ 7.49 :
\]
sembilan bilangan, satu untuk tiap tanda skala yang baru.

\textbf{3.} Gambar yang sama, sepuluh kali lebih kecil: antara $7.440$
dan $7.450$ terletak $7.441, 7.442, \dots, 7.449$ --- sembilan
bilangan lagi.

\textbf{4.} Antara $7.400$ dan $7.500$, bilangan dengan tiga angka
desimal adalah $7.401, 7.402, \dots, 7.499$: semua perseribuan dari
$401$ sampai $499$, yaitu $99$ bilangan.

\textbf{5.} Resep perbesarannya: tulislah kedua bilangan dengan
banyak tempat desimal yang sama, lalu \emph{tambahkan satu tempat
desimal lagi} --- bilangan yang lebih kecil diikuti sebuah angka dari
$1$ sampai $9$ mendarat tepat di antara keduanya. Memang
$5.67 = 5.670$ dan $5.68 = 5.680$, dan
\[
5.670 < 5.675 < 5.680 ;
\]
begitu pula $0.1999 = 0.19990 < 0.19995 < 0.20000 = 0.2$. Di antara
dua tanda skala yang bertetangga selalu ada satu tingkat baru berisi
sembilan skala yang lebih halus.

\textbf{6.} $3 < 3.05 < 3.1$; lalu $3 < 3.005 < 3.01$; lalu
$3 < 3.0005 < 3.001$. Setiap calon dikalahkan oleh satu perbesaran lagi.

\textbf{7.} Bilangan apa pun yang diajukan Tono --- sebut saja
calonnya, yang lebih besar daripada $3$ --- pertanyaan 5
menghasilkan bilangan yang terletak tepat di antara $3$ dan calon
itu. Jadi calonnya \emph{bukan} bilangan yang paling dekat dengan
$3$: ada yang lebih dekat lagi. Karena hal ini terjadi pada setiap
calon tanpa kecuali, tidak ada bilangan yang dapat menjadi
``bilangan tepat sesudah $3$'': bilangan itu memang tidak ada.

\textbf{8.} $2.9 < 2.95 < 3$, \ $2.99 < 2.995 < 3$, \
$2.999 < 2.9995 < 3$. Dan $3 - 2.999 = 0.001$. Dengan dua puluh
angka sembilan, \omterm{ex:g1:subtraction:difference}{selisihnya} adalah sebuah \omterm{def:g4:decimals:point}{titik desimal} yang diikuti
sembilan belas \omterm{def:g1:counting:zero}{nol} dan sebuah $1$ --- satu satuan pada tempat desimal
kedua puluh: sangat kecil, tetapi bukan \omterm{def:g1:counting:zero}{nol}. Sebanyak apa pun angka
sembilan yang ditulis Tono, ia tidak pernah mencapai $3$.

\textbf{9.} Bilangan cacah menaiki tangga selangkah $1$: lebih dari
$7$ tetapi kurang dari $8$ berarti bilangan cacah satuan yang
terletak tepat di antara $7$ dan $8$ satuan, dan bilangan itu tidak
ada. Perbesaran meloloskan diri dari hal ini hanya dengan menulis
angka \emph{sesudah titik} --- justru yang tidak dimiliki bilangan cacah.

\textbf{10.} Panjang terkecil yang membulat menjadi $3.7$ adalah
$3.65$ cm: angka berikutnya $5$, yang \omterm{def:g4:numbers:round}{dibulatkan} \emph{ke atas}
(\cref{def:g6:decimals:rounding}). Dan $3.75$ membulat naik menjadi
$3.8$ karena alasan yang sama. Jadi ``$3.7$ cm ke persepuluhan
terdekat'' berarti: sedikitnya $3.65$ cm, dan tepat di bawah $3.75$
cm --- seluruh ruas yang diperbesar berisi panjang yang mungkin,
bersembunyi di balik satu nilai yang \omterm{def:g4:numbers:round}{dibulatkan}.

\textbf{11.} $0.0999 < 0.123 < 0.58 < 0.6$. Satu kalimat:
membandingkan angka demi angka dari kiri, $0.0999$ punya $0$
persepuluhan sedangkan $0.6$ punya $6$, dan perbandingannya selesai
di situ --- banyaknya angka yang ditulis tidak mengatakan apa-apa
tentang besarnya (\cref{met:g6:decimals:compare}).

\textbf{12.} Tidak: dalam sen, $4.99$ euro adalah $499$ sen dan $5$
euro adalah $500$ sen --- bilangan \emph{cacah} yang berurutan, tanpa
apa pun di antaranya. Harga, yang punya tepat dua tempat desimal,
sebenarnya adalah bilangan cacah sen yang menyamar: seperti bilangan
cacah pada pertanyaan 9, harga memang punya tetangga sebelah.
Perbesaran tanpa henti memerlukan hak untuk menulis tempat desimal
yang makin banyak.

\textbf{13.} Calon di dekat $6$ adalah $5.82$ dan $8.25$ (bilangan
yang diawali $2$, atau $58$, $25$, $82$, $85$, jauh dari $6$).
Jaraknya: $6 - 5.82 = 0.18$ dan $8.25 - 6 = 2.25$. Yang paling dekat adalah $5.82$.

\textbf{14.} Misalnya $3.1415$, karena
$3.1410 < 3.1415 < 3.1420$; lalu $3.14159$, karena
$3.14150 < 3.14159 < 3.14160$. (Keduanya adalah langkah nyata dalam
perburuan $\pi = 3.14159\dots$ yang sesungguhnya.)

\textbf{15.} Andaikan seseorang menyebut sebuah bilangan, sebesar apa
pun yang ia mau. Satu perbesaran mengubah setiap pasang skala
bertetangga antara $0$ dan $1$ menjadi sembilan bilangan baru, dan
perbesaran itu dapat diulang selamanya --- pertanyaan 2, 3 dan 4
hanyalah dua tingkat yang pertama. Mengulangnya cukup banyak kali
menghasilkan \omterm{def:g6:decimals:places}{bilangan desimal} antara $0$ dan $1$ yang lebih banyak
daripada bilangan yang disebut tadi. Jadi tidak ada bilangan yang
cukup besar untuk membilangnya: banyaknya tak hingga.

\end{solution}
